\chapter{Software Implementation and Experimental Setup}

\section{Implementation Structure}
The repository is organized into a reusable Python package and command-line
experiments. The Phase 1 data and model modules are kept separate from the
Phase 2 communication primitives. This separation allows the baseline to be
retrained independently and prevents the communication experiment from
depending on a hidden or manually edited result.

\begin{table}[H]
\centering
\caption{Software modules used in the project.}
\begin{tabular}{p{5cm} p{8cm}}
\hline
\textbf{File} & \textbf{Function}\\
\hline
\texttt{src/amc/data.py} & RadioML loading, normalization, stratified splits\\
\texttt{src/amc/model.py} & Phase 1 Conv1D CNN definition\\
\texttt{src/amc/communications.py} & Modulation, channel, estimation, demodulation\\
\texttt{scripts/train.py} & Phase 1 model training and checkpointing\\
\texttt{scripts/evaluate.py} & Phase 1 test and SNR-wise evaluation\\
\texttt{scripts/run\_phase2\_link.py} & Complete Phase 2 link simulation\\
\texttt{scripts/build\_final\_report.py} & Consolidated comparison figure\\
\texttt{scripts/build\_parameter\_figures.py} & Parameter and architecture figures\\
\hline
\end{tabular}
\end{table}

\section{Phase 1 Dataset Preparation}
RadioML 2016.10a is loaded from its dictionary representation. Each key
contains a modulation name and SNR value, while the associated value contains
I/Q arrays. The loader concatenates all groups, assigns integer class labels,
records SNR values, and performs per-example power normalization. A
deterministic stratified split is then generated using seed 42.

\section{Phase 1 Model Configuration}
The baseline model receives two input channels and 128 temporal samples. The
first convolution uses 64 filters with kernel size 7, the second uses 128
filters with kernel size 5, and the third uses 256 filters with kernel size 3.
Each convolution is followed by batch normalization and ReLU. Max-pooling
reduces the temporal dimension, adaptive average pooling creates a fixed-size
representation, and dropout reduces overfitting. The output layer contains
11 class logits and the complete network has 144,331 trainable parameters.

\section{Training Procedure}
Training uses class-weighted cross-entropy, AdamW optimization, cosine
annealing, a learning rate of $3\times10^{-4}$, weight decay $10^{-4}$,
batch size 1024, and seed 42. The reported run uses eight epochs. The best
validation checkpoint is stored with class ordering, input shape, validation
history, and model weights.

\section{Phase 2 Link Configuration}
The final Phase 2 experiment uses BPSK, QPSK, 8-PSK, and 16-QAM with 256
symbols per frame. The SNR is swept from -4 dB to +16 dB in 2 dB increments,
with 200 frames per SNR and seed 42. Every frame passes through block
Rayleigh fading, AWGN, phase offset, and frequency offset.

\begin{table}[H]
\centering
\caption{Final experimental settings.}
\begin{tabular}{p{6cm} p{7cm}}
\hline
\textbf{Parameter} & \textbf{Value}\\
\hline
Modulation formats & BPSK, QPSK, 8-PSK, 16-QAM\\
Frame size & 256 symbols\\
Channel & Block Rayleigh fading plus AWGN\\
Carrier impairments & Random phase and frequency offset\\
Synchronization & Pilot-assisted simulated estimate\\
SNR sweep & -4, -2, \ldots, 16 dB\\
Frames per SNR & 200\\
Metrics & BER, SER, packet success, throughput, decision accuracy\\
Seed & 42\\
\hline
\end{tabular}
\end{table}

\section{Execution Commands}
The Phase 2 experiment is executed from the repository root:
\begin{verbatim}
python scripts/run_phase2_link.py --frames 200 \
  --output artifacts/phase2-link-results.png
\end{verbatim}
The command writes both a graph and a JSON file. The JSON file is the source
for the final comparison and parameter figures:
\begin{verbatim}
python scripts/build_final_report.py
python scripts/build_parameter_figures.py
\end{verbatim}

\section{Verification}
Four automated tests verify reproducible data splitting, input validation,
modulation round trips, and clean constellation identification. Python
compilation and whitespace checks are also run before publication.
